\documentclass[10pt,a4paper]{amsart}
\usepackage[margin=19mm]{geometry}
\usepackage{amsmath,amssymb,mathtools}
\usepackage[scr=rsfs,cal=euler]{mathalfa}
\usepackage{xfrac}
\usepackage[unicode,hidelinks,bookmarks=false]{hyperref}
\newcommand{\ie}{i.e.\ }
\newcommand{\cf}{cf.\ }
\newcommand{\resp}{respectively\ }
\newcommand{\Subsection}{\subsection}
\providecommand{\nobreakdash}{\nobreak\hbox{-}}
\setlength{\emergencystretch}{2em}
\multlinegap0pt
\usepackage{bm}
\usepackage{bbm}
\usepackage{dsfont}
\usepackage{xspace}
\usepackage{cancel}
\usepackage{mathtools}
\usepackage{amsthm}
\makeatletter
\@ifundefined{theorem}{\newtheorem{theorem}{Theorem}}{}
\makeatother
\newcommand{\E}{\mathrm E}
\newcommand{\HH}{\mathrm H}
\newcommand{\LL}{\mathrm L}
\newcommand{\ind}{\mathbf1}
\title[Cayley-graph density of Thompson's group F: local deletions ]{Cayley-graph density of Thompson's group F: local deletions and finite-window bounds}
\author{Thomas Prellberg}
\date{Source: arXiv:2609.12290v2 (CC BY 4.0)}
\begin{document}
\maketitle

\noindent\emph{Source and licence.} Cayley-graph density of Thompson's group F: local deletions and finite-window bounds. Version arXiv:2609.12290v2 (CC BY 4.0), distributed under the Creative Commons Attribution 4.0 International licence, as recorded in the arXiv licence metadata for this exact version and shown on the abstract page https://arxiv.org/abs/2609.12290v2. Full rights notice: licenses/arxiv-2609.12290-CC-BY-4.0.txt.\par\smallskip

\noindent\textit{Editorial scope.} Counted unit: the Theorem whose source label is \texttt{thm:one-right} in the article (source0.tex lines 821--823, source label \texttt{thm:one-right}), delivered together with the article's own complete original proof of it (source0.tex lines 824--864, one \texttt{proof} environment of 41 lines). No mathematics is written, repaired or re-derived: statement and proof are the article's own text line for line. Required context retained uncounted: eq:allocation (lines 799--803). Every definition, display and label of those lines is retained unchanged. Omitted from this package and not redistributed: 1--798, 804--820, 865--918. Nothing inside a retained excerpt is omitted or altered: every retained body is delivered exactly as the article's lines are written, so no omission receipt of any kind accompanies this package. Citations: the retained excerpts carry no citation command at all, so no bibliography entry of the article is transcribed and no reference is redistributed. Reference adaptation: none was needed and none was made; every reference inside the delivered unit resolves inside it, so no unresolved reference or citation is produced. Printed slips kept literal: no printed word, symbol, hypothesis or number of the counted statement or of its proof was altered. Layout and class: the article's own class is \texttt{article} with packages including fontenc, lmodern, amsmath, amssymb, amsthm, mathtools, booktabs, array, tikz, geometry, microtype, hyperref, enumitem, needspace; the wrapper is the lane's pinned \texttt{amsart} editorial wrapper at 10pt on A4, which restates the theorem-like environments the retained text needs and adds the article's own macro definitions that the retained text needs (\texttt{\textbackslash E}, \texttt{\textbackslash HH}, \texttt{\textbackslash LL}, \texttt{\textbackslash ind}), with extra wrapper packages (bm, bbm, dsfont, xspace, cancel, mathtools). The delivered numbering is the unit's own. Assets: no retained span contains a figure environment, a graphics-inclusion command, a PDF-page inclusion, a table or a TikZ picture; no image file is copied into this package, the package declares no asset dependency and no image file is redistributed with it. Licence: version arXiv:2609.12290v2 of this article is distributed under the Creative Commons Attribution 4.0 International licence, as recorded in the arXiv licence metadata for this exact version and shown on the abstract page https://arxiv.org/abs/2609.12290v2. A full rights notice is delivered at licenses/arxiv-2609.12290-CC-BY-4.0.txt. Characters: every retained excerpt is pure ASCII, so the delivered engine, XeLaTeX with its default Latin Modern text font, reports no missing character.
\begin{align}\label{eq:allocation}
0\le A_{uv}&\le2m_{uv},\notag\\
2m_{uu}+\sum_{v>u}A_{uv}+\sum_{v<u}(2m_{vu}-A_{vu})&\le d\mu_u
\quad\text{for every }u.
\end{align}

\begin{theorem}[One right neighbour]\label{thm:one-right}
For every fixed $a\ge0$, among all coarse-category retention rules on the window $[-a,1]$, the maximum limiting density is $7/2$. The limit is taken first as $n\to\infty$ and then as $k\to\infty$, as above.
\end{theorem}

\begin{proof}
Set
\[
z=u_2=\sqrt\ell,\qquad \theta=\frac{z-e}{\ell},\qquad
(t_\E,t_\LL,t_\HH)=(1,\theta,0),\qquad \eta_b=\ind_{b=\HH}.
\]
Thus $\sum_b p_bt_b=z$, $\sum_b p_b\eta_b=h$, and $h+z^2=3/4$.
Define a horizontal bias, indexed by two categories, by
\begin{equation}\label{eq:onebias}
f_{bc}=2\eta_b+2\eta_c+2t_c(t_b+z)-3.
\end{equation}
In the order $\E,\LL,\HH$ its nine entries are
\begin{equation}\label{eq:biasmatrix}
(f_{bc})=
\begin{pmatrix}
2z-1&2\theta(1+z)-3&-1\\
2(\theta+z)-3&2\theta(\theta+z)-3&-1\\
2z-1&2\theta z-1&1
\end{pmatrix}.
\end{equation}
All lie in $[-1,1]$: the elementary bounds $3/4<z<4/5$ and $4/5<\theta<9/10$ suffice for each displayed entry.

We allocate the doubled mass of every elementary edge before aggregating equal endpoint types. On a rightward marker edge, let $c,d$ be the marked category and its right neighbour \emph{at the target}. Assign $1+f_{cd}$ times the edge mass to its source, and $1-f_{cd}$ times the mass to its target. For a split with parent in $\HH$, assign all doubled mass to the parent endpoint. For a split with parent in $\LL$, assign it all to the left-child endpoint. Every allocation is nonnegative and uses exactly the doubled edge mass.

Consider a type ending in $b,c$, of mass $\mu$. The total allocated split mass at this type, divided by $\mu$, is
\[
2\eta_b+2t_bt_c.
\]
The first term is the contribution of its own height-$k$ root; the second is the joint mass $\kappa_{\LL bc}=(v_2)_b(v_2)_c$, divided by $p_bp_c$. The $a$ preceding categories have the same product weight on both split endpoints and cancel from this quotient.

Put $g_c=\sum_d p_df_{cd}$. Equation~\eqref{eq:onebias} and $h+z^2=3/4$ give
\[
g_c=2\eta_c+2zt_c-\frac32,\qquad
f_{bc}=2\eta_b+2t_bt_c-\frac32+g_c.
\]
The rightward edge allocation contributes $1+g_c$ per unit vertex mass; incoming rightward edges contribute $1-f_{bc}$. Including the split allocation, the total is consequently
\[
2+g_c-f_{bc}+2\eta_b+2t_bt_c=\frac72.
\]
This calculation is independent of $a$. Summing the nonnegative allocations over any retained collection of types proves the upper bound as in \eqref{eq:allocation}. Retaining all types attains $7/2$, proving equality.
\end{proof}

\end{document}
